% As micrografias só são incluídas quando os arquivos científicos estão locais.
% Não usar ilustrações sintéticas ou fotografias de encapsulamento como substitutas.
\IfFileExists{assets/microscopia/vlsi_nise.jpg}{%
\begin{frame}{Dentro do chip: muitas conexões, muitos dispositivos}
\begin{center}\includegraphics[width=.92\linewidth,height=4.5cm,keepaspectratio]{assets/microscopia/vlsi_nise.jpg}\end{center}
\chamada{A organização do circuito aparece ao ampliar suas camadas.}
\fonte{NISE Network / Lawrence Hall of Science, \href{https://www.nisenet.org/catalog/media/zoom_microchip_video}{Zoom into a Microchip}, recurso de 2013. \href{https://creativecommons.org/licenses/by-nc-sa/3.0/us/}{CC BY-NC-SA 3.0 US}. Imagem estática; circuito mais antigo e processo não identificado.}
\end{frame}}{}
\IfFileExists{assets/microscopia/gaa_imec.jpg}{%
\begin{frame}{Do circuito ao transistor nanométrico: uma vista TEM}
\begin{center}\includegraphics[width=.92\linewidth,height=4.5cm,keepaspectratio]{assets/microscopia/gaa_imec.jpg}\end{center}
\chamada{Nanosheets: canais empilhados envolvidos pelo Gate.}
\fonte{imec, \href{https://www.imec-int.com/en/articles/outer-wall-forksheet-bridge-nanosheet-and-cfet-device-architectures-logic-technology}{Outer wall forksheet: bridging nanosheet and CFET}, Fig. 1 (2025). Microscopia eletrônica de transmissão (TEM). Amostra diferente da vista NISE; não é um zoom contínuo do mesmo chip.}
\end{frame}}{}
